\begin{satz}{Fortsetzung einer endlichen Erwartungswertungleichung (II)}
            {AbzaehlbareDarstellungUndSupremumVonSuprema}
Seien $I, J$ nichtleere, abzählbare Mengen,
$X, Y: I\times J \rightarrow \mcm(\Omega, \mfa, \R)$, sowie
$\mci := \Pot_{ne}(I)$ und $\mcj := \Pot_{ne}(J)$.\\
Es gelte:
\begin{enumerate}[(a)]
\item $\forall\,\omega\in\Omega\,\forall\, Z\in \meng{X, Y}: 
\menge{Z(i, j)(\omega)}{(i, j)\in I\times J}$
      nach oben beschränkt.
\item $\forall A \in \mci \cup \meng I \forall B \in \mcj \cup \meng J
\forall Z\in \meng{X, Y}:
     \underset{i \in A}\sup\,\underset{j \in B}\sup\, Z(i, j)
     \in L_1(\Omega, \mfa, \Prim, \R)$.
\item $\forall\, A \in \mci \, \forall\, B \in \mcj: 
\E\bigg(\underset{i \in A}\sup\,\underset{j \in B}\sup\, Y(i, j) \bigg)
\le \E\bigg(\underset{i \in A}\sup\,\underset{j \in B}\sup\, X(i, j) \bigg)$.\\
 (Alle genannten Abbildungen sind wegen $(a)$ und
$I \neq \emptyset \neq J$wohldefiniert.)
\end{enumerate}
Dann gilt
\[\E\left(\sup_{i \in I}\sup_{j \in J}Y(i, j) \right) \le
  \E\left(\sup_{i \in I}\sup_{j \in J}X(i, j) \right)\text.\]
\end{satz}
\begin{proof}
Seien $\mcp:= \Pot_{ne}(I \times J)$ und $(i_0, j_0) \in I \times J$.\\
Dann sind $Y(i_0, j_0)$ und
\[\sup_{(i, j) \in I \times J}Y(i, j)= \sup_{i \in I} \sup_{j \in J}Y(i, j)
\text{ nach (b) integrierbar.}\]
Da $I$ und $J$ abzählbar sind, ist auch $I \times J$ abzählbar.
Zusammen mit (a) sind damit die Voraussetzungen von
\ref{EndlicheDarstellungDesSupremumsIntegrierbarerFunktionen} für die Familie
$(Y(i,j))_{(i,j) \in I \times J}$ erfüllt. Also ist
\[M_1:=\menge{\E\left(\sup_{(i, j) \in P}Y(i,j) \right)}{P \in \mcp}\]
nach \refclaim{EndlicheDarstellungDesSupremumsIntegrierbarerFunktionen}
{EndlicheDarstellungDesSupremumsIntegrierbarerFunktionen1} nichtleer und nach
oben beschränkt und es gilt
\begin{align}\sup M_1
\overset{\refclaim{EndlicheDarstellungDesSupremumsIntegrierbarerFunktionen}
{EndlicheDarstellungDesSupremumsIntegrierbarerFunktionen2}}= 
\E\bigg(\sup_{(i, j) \in I \times J}Y(i, j) \bigg)
= \E\left(\sup_{i \in I}\sup_{j \in J}Y(i,j)\right)\text. \label{feeeiieq:1}
\end{align}
Sei nun 
\[M_2 := \menge{\E\left(\sup_{i \in A}\sup_{j \in B}Y(i,j) \right)}
               {A \in \mci, B \in \mcj}\text.\]
Offenbar ist $M_2 \subseteq M_1$ und wegen $I \neq \emptyset\neq J$ auch
nichtleer.
Da $M_1$ nach oben beschränkt ist, ist es auch
$M_2$ und es gilt $\sup M_2 \le \sup M_1$.\\
Sei nun $Q \in \mcp$. Dann gibt es ein
$A_Q \in \mci$ und ein $B_Q \in \mcj$ mit 
$Q \subseteq A_Q \times B_Q$.\\
Da $Q \neq \emptyset$ existiert ein $(i', j') \in Q$, sodass gilt:
\[Y(i',j')\le\sup_{(i,j)\in Q}Y(i,j)\le\sup_{(i,j)\in A_Q\times B_Q}Y(i,j)
=\sup_{i\in A_Q}\sup_{j\in B_Q}Y(i,j)\text.\]
Nach (b) sind die erste und die letzte Abbildung integrierbar, also ist
es nach \refclaim{Integrierbarkeitskriterien}{Integrierbarkeitskriterien2}
auch die zweite Abbildung. Somit gilt:
\[\E\left(\sup_{(i, j) \in Q}Y(i, j) \right) \le 
\E\left(\sup_{(i, j) \in A_Q \times B_Q}Y(i, j) \right)
= \E\left(\sup_{i \in A_Q}\sup_{j \in B_Q}Y(i, j) \right) \in M_2\text.\]
Erhalte: $\forall \, x \in M_1\, \exists \, y \in M_2: x \le y$.\\
Damit ist offensichtlich:
\begin{align}
 \sup M_1\le\sup M_2\text{, also }
\sup M_1=\sup M_2\text. \label{feeeiieq:2}
\end{align}
Seien nun $C \in \mci$ und $D \in \mcj$. Dann gilt
\[\sup_{i\in C}\sup_{j\in D}X(i,j)\le\sup_{i\in I}\sup_{j \in J}X(i,j)\text.\]
Gemäß (b) sind beide Abbildungen integrierbar und daher ist auch
\[\E\left(\sup_{i \in C} \sup_{j \in D}X(i, j) \right)
\le \E\left(\sup_{i \in I} \sup_{j \in J}X(i, j) \right)\text.\]
Daher gilt
\[\E\left(\sup_{i \in C} \sup_{j \in D}Y(i, j) \right)
\overset{(c)}\le\E\left(\sup_{i \in C} \sup_{j \in D}X(i, j) \right)
\le \E\left(\sup_{i \in I} \sup_{j \in J}X(i, j) \right)\]
und damit
\begin{align}
\sup\menge{\E\left(\sup_{i \in A} \sup_{j \in B}Y(i, j) \right)}
            {A \in \mci, B \in \mcj} \le
            \E\left(\sup_{i \in I} \sup_{j \in J}X(i, j) \right)\text.
            \label{feeeiieq:3}
\end{align}
Schließlich gilt also
\begin{align*}
&\,\E\left(\sup_{i \in I}\sup_{j \in J}Y(i,j)\right)
\overset{\eqref{feeeiieq:1}}=
\sup M_1 \overset{\eqref{feeeiieq:2}}= \sup M_2\\
= &\, \sup \menge{\E\left(\sup_{i \in A}\sup_{j \in B}Y(i,j) \right)}
        {A \in \mci, B \in \mcj}
 \overset{\eqref{feeeiieq:3}}\le
 \E\left(\sup_{i\in I}\sup_{j \in J}X(i, j) \right)\text,
\end{align*}
was gerade behauptet wurde.
\end{proof}